% from ega4-2-fr.tex
\oldpage[IV-2]{61}

relation d'équivalence $R$ définie par $p$ sur $X'$ est ouverte, donc il en est
de même de la relation $R_{Z'}$ qu'elle induit sur la partie saturée $Z'$, et
$Z$ s'identifie à l'espace quotient $Z'/R_{Z'}$ (Bourbaki,
\emph{Top. gén.}, chap.~I\textsuperscript{er}, 3\textsuperscript{e} éd.,
\S~5, n\textsuperscript{o}~2, prop.~4)~; on peut donc appliquer à $Z$ et $Z'$
le lemme (4.4.2), d'où la première assertion de (4.4.4).
\end{proof}

\begin{corollary}[4.4.6]
\label{IV.2.4.4.6-fr}
Les hypothèses et notations étant celles de (4.4.5), l'application
$Z\mapsto p^{-1}(Z)$ est une bijection de l'ensemble des composantes
irréductibles (resp. connexes) de $X$ sur l'ensemble des composantes
irréductibles (resp. connexes) de $X\otimes_k K$, dont la bijection réciproque
est $Z'\mapsto p(Z')$.
\end{corollary}

\subsection{Préschémas géométriquement irréductibles et géométriquement
connexes}
\label{subsection:IV.2.4.5-fr}

\begin{proposition}[4.5.1]
\label{IV.2.4.5.1-fr}
Soient $k$ un corps, $X$ un $k$-préschéma, $\Omega$ une extension
algébriquement close de $k$. Soit $n$ (resp. $n'$) le cardinal de l'ensemble
des composantes irréductibles (resp. des composantes connexes) de
$X\otimes_k\Omega$~; ce nombre est indépendant de l'extension algébriquement
close $\Omega$ choisie. En outre, pour toute extension $K$ de $k$, le cardinal
de l'ensemble des composantes irréductibles (resp. des composantes connexes)
de $X\otimes_k K$ est $\leqslant n$ (resp. $\leqslant n'$).
\end{proposition}

\begin{proof}
En effet, deux extensions algébriquement closes $\Omega$, $\Omega'$ de $k$
peuvent toujours être considérées comme des sous-extensions d'une même
extension $E$ de $k$~; il suffit alors d'appliquer (4.4.6) à
$X\otimes_k\Omega$ et $X\otimes_k\Omega'$ et à l'extension commune $E$ de
$\Omega$ et $\Omega'$ pour démontrer la première assertion. La seconde
s'obtient en prenant pour $\Omega$ une extension algébriquement close de $K$
et en utilisant (4.4.1).
\end{proof}

\begin{definition}[4.5.2]
\label{IV.2.4.5.2-fr}
Le cardinal $n$ (resp. $n'$) de (4.5.1) s'appelle le \emph{nombre géométrique
de composantes irréductibles} (resp. \emph{de composantes connexes}) de $X$
(relativement à $k$). Si $n=1$ (resp. $n'\leqslant1$) on dit que $X$ est un
$k$-préschéma \emph{géométriquement irréductible} (resp. \emph{géométriquement
connexe}).
\end{definition}

On retrouve ainsi une définition donnée pour un cas particulier dans
(III, 4.3.4). En vertu de (4.5.1), les propriétés suivantes sont
équivalentes~:
\begin{enumerate}
\item[$a)$] $X$ est géométriquement irréductible (resp. géométriquement
connexe).
\item[$b)$] Pour une extension algébriquement close $\Omega$ de $k$,
$X\otimes_k\Omega$ est irréductible (resp. connexe).
\item[$c)$] Pour toute extension $K$ de $k$, $X\otimes_k K$ est irréductible
(resp. connexe).
\end{enumerate}

\begin{env}[4.5.3]
\label{IV.2.4.5.3-fr}
Soient $X$ un $k$-préschéma, $Z$ une partie localement fermée de $X$. Si l'on
désigne encore par $Z$ un quelconque des sous-préschémas de $X$ ayant $Z$ pour
espace sous-jacent (I, 5.2.1), il résulte de (I, 5.1.8) que pour une extension
$K$ de $k$, le nombre de composantes irréductibles (resp. connexes) de
$Z\otimes_k K$ est indépendant du sous-préschéma d'espace sous-jacent $Z$ que
l'on a choisi~; aussi définit-on le nombre géométrique de composantes
irréductibles (resp. de composantes connexes) de $Z$ comme celui de tout
sous-préschéma de $X$ ayant $Z$ pour espace sous-jacent.
\end{env}

\begin{proposition}[4.5.4]
\label{IV.2.4.5.4-fr}
Soient $X$, $Y$ deux $k$-préschémas, $f:X\to Y$ un $k$-morphisme surjectif.
Si $X$ est géométriquement irréductible (resp. géométriquement connexe), il en
est de même de $Y$.
\end{proposition}

